\section{Unordered selections: combinations}

A combination is an unordered selection.

\begin{definitionbox}
A \textbf{combination} of size \(k\) from an \(n\)-element set is a subset with \(k\) elements.
\end{definitionbox}

The number of combinations is
\[
\binom{n}{k}=\frac{n!}{k!(n-k)!}.
\]

This formula can be derived from ordered selections. There are
\[
P(n,k)=\frac{n!}{(n-k)!}
\]
ordered selections of length \(k\). But each unordered selection of \(k\) objects can be ordered in \(k!\) ways. Therefore
\[
\binom{n}{k}=\frac{P(n,k)}{k!}=\frac{n!}{k!(n-k)!}.
\]

Combinations are used when:
\begin{itemize}
    \item only membership is recorded,
    \item order is ignored,
    \item each object may be chosen at most once,
    \item the outcome is a set.
\end{itemize}

\begin{examplebox}
A five-card hand from a standard deck is a set of five cards, not a sequence of five draws, if the final hand is all that is recorded. The number of possible hands is
\[
\binom{52}{5}.
\]
If the order of drawing is recorded, the correct count is instead
\[
P(52,5)=52\cdot 51\cdot 50\cdot 49\cdot 48.
\]
These are different sample spaces.
\end{examplebox}
